\chapter{世界如何进入大脑}
% 完整版：chapters/11_sensors.tex

任何感觉器官都是一个换能器，就像麦克风或光电管。光、声音或分子作用于传感细胞里的一种特殊蛋白质，这种蛋白质打开一道“阀门”，电流流过，到链条末端便产生一串咔嗒，送进网络。几乎所有传感器都释放谷氨酸，所以机器的输入直接接到了它的电话网上。

\section*{在物理极限上}

大脑的传感器工作在可能性的极限边缘。在黑暗中，光传感器能对单个光粒子作出响应。声音传感器能察觉不到一纳米的位移——比几个原子还小。在昏暗中，限制精度的已经不是眼睛，而是光本身：光粒子随机到达，要分辨色调，就得多攒一些。所以黄昏时我们看东西更慢——眼睛攒光的时间更长，就像长曝光的相机。

\section*{自动曝光}

从星光下的夜晚到阳光下的雪地，照度相差几十亿倍，声音的响度差得更多。而咔嗒的频率只在零到每秒几百次之间变化。把前者直接塞进后者是不可能的。解决办法和相机一样：自动曝光。传感器不断根据周围的平均水平调整灵敏度，报告的不是亮度，而是\emph{相对于背景}的亮度。不变的东西久而久之不再引起响应，而每一次变化都会。机器的输入从一开始说的就是变化。

\pencilfig{125}{%
% light rises in two tenfold steps, then drops a hundredfold; sensor rate = baseline + gain * (log light - adapted level),
% the adapted level follows log light with a 0.7 s time constant; rate clipped at zero
\draw[penlite=pcAmber] (0,1.75) -- (1.4,1.75) -- (1.4,2.1) -- (4.2,2.1) -- (4.2,2.45) -- (6.3,2.45) -- (6.3,1.75) -- (8.4,1.75);
\node[lbl, text=pcAmber!70!black, anchor=east] at (-0.05,2.0) {光};
\node[lbl, anchor=south] at (2.8,2.12) {亮10倍}; \node[lbl, anchor=south] at (5.25,2.47) {再亮10倍};
\node[lbl, anchor=south] at (7.35,1.77) {又暗了};
\draw[pen] (0,0) -- (8.4,0);
\draw[pcGray, densely dashed, line width=0.5pt] (0,0.32) -- (8.4,0.32);
\draw[penlite=pcBlue] plot coordinates {(0.00,0.32) (1.26,0.32) (1.40,1.02) (1.54,0.85) (1.68,0.72) (1.82,0.62) (1.96,0.54) (2.10,0.49) (2.24,0.45) (2.38,0.41) (2.52,0.39) (2.66,0.37) (2.80,0.36) (3.08,0.34) (3.50,0.33) (4.06,0.32) (4.20,1.03) (4.34,0.85) (4.48,0.72) (4.62,0.62) (4.76,0.54) (4.90,0.49) (5.04,0.45) (5.18,0.41) (5.32,0.39) (5.46,0.37) (5.60,0.36) (5.88,0.34) (6.16,0.33) (6.30,0.00) (7.00,0.00) (7.14,0.07) (7.28,0.13) (7.42,0.18) (7.56,0.22) (7.70,0.24) (7.84,0.26) (7.98,0.28) (8.12,0.29) (8.40,0.30)};
\node[lbl, text=pcBlue, anchor=east] at (-0.05,0.32) {响应};
\node[lbl, anchor=north] at (6.65,-0.02) {沉默};
}{光一级一级变亮，每级十倍。传感器对每次变化都报以一次爆发，然后很快回到原来的水平。重新变暗时，它会沉默一阵子。}

工程师在事件相机里复制了这一招。这种相机不按固定节拍拍摄画面：每个像素在有变化时自己报告“变亮了”或“变暗了”。这正是第二章那对“是”和“否”导线，只不过做在了硅片上。

\section*{眼睛压缩图像}

眼睛里大约有一亿个光传感器，而把图像送往大脑的导线只有大约一百万根。图像在离开眼睛之前，要被压缩约一百倍。平坦均匀的区域几乎不占带宽，传出去的是边缘和变化。根据在动物眼睛上做的估算、再换算到人身上，眼睛向大脑传送的信息量约为每秒一千万比特——相当于一根老式网线。

\section*{耳朵是一架钢琴}

耳朵把声音按频率分解开，就像工程师会装上一组滤波器那样。耳朵里有一块弹性隔膜，不同位置对不同频率起反应，每一段上的传感器只听自己的那个“琴键”。被触发的传感器的编号，就是音高。

\pencilfig{11}{\pencilart{11}}{耳朵像一架反过来的钢琴：每个琴键只听自己的那个音。}

不仅如此，耳朵还不是被动的。它的一部分细胞会随着声音的节拍主动推动隔膜，把轻声放大几千倍，而对响声几乎不放大。这是一个工作在自激边缘的放大器：在那里它格外灵敏。在低频时，神经元还会跟着声波的节拍咔嗒，保留声音的精确时刻——大脑正是靠它来确定声源的方向。频率越高，神经元越跟不上，最后只剩下位置编码。

\section*{其他感觉}

几乎在每一种感觉里，都能认出前面几章的某个手法。温度由两根导线传递，冷和热各一根。触觉既用一碰就响应、随即沉默的传感器，也用只要压力还在就一直保持响应的传感器。而嗅觉像条形码：人有大约四百种嗅觉传感器，一种气味就是一组被触发的类型。这样的组合数目是个天文数字。

\fullversion{11}
